\documentclass{subfiles}
\begin{document}
    The first kind of optimization algorithms we consider are the so-called \emph{greedy algorithms}.
        Although these differ for the problem they solve,
        all greedy algorithms must follow the same structure, which we summarize below.
        \begin{enumerate}
            \item Let $X$ be the domain of our problem.
                The algorithms must maintain a partition $\set{Y, W, Z}$ of $X$,
                where $Y$ is the set of elements that are part of the final solution,
                $W$ is the set of discarded elements and $Z$ that of elements yet to examine.

            \item There cannot be backtracking; that is, 
                we can't go back and modify the current solution,
                i.e., we can remove elements from $Y$ and add some element of $W$.

            \item The elements from $Z$ are choosed according to some greedy mechanism.
                Meaning that an element is accepted only if it satisfies some local optimality criteria.
        \end{enumerate}
        Moreover, in general, the number of steps of a greedy algorithm is known.

    \subsection{No Free Lunch Theorem}
    \subfile{../sub/2.1 - NFLT}

\end{document}
